\documentclass[11pt]{article}
\usepackage[margin=1in]{geometry}
\usepackage{amsmath,amssymb}
\usepackage{booktabs}
\usepackage{array}
\usepackage{xcolor}
\usepackage{listings}
\usepackage[hidelinks]{hyperref}
\usepackage{microtype}

\newcommand{\todo}[1]{\textcolor{red}{[\textbf{TODO:} #1]}}
\newcommand{\hx}[1]{\texttt{#1}}
\lstset{basicstyle=\ttfamily\footnotesize,breaklines=true,frame=single,columns=fullflexible}

\title{A Practical Collision for 36-Step SHA-512\\[2pt]\large (Draft --- not for distribution)}
\author{Junhong Huang\\ Boston University \quad \href{mailto:raymxnd@bu.edu}{\texttt{raymxnd@bu.edu}}}
\date{22 November 2026}

\begin{document}
\maketitle

\begin{abstract}
We give a colliding message pair for SHA-512 reduced to 36 steps, starting from the standard initial value. To our knowledge it is the first practical collision for 36-step SHA-512; the previous practical record was 35 steps~\cite{LLWS26}. The attack is not new. It uses the 36-step differential characteristic and the memory-efficient two-block framework of Li~et~al.~\cite{LZLLQZ26}, whose estimated time complexity is $2^{66.2}$. Our contribution is to make it practical by:
(i)~counting the degrees of freedom and the uncontrolled conditions more precisely;
(ii)~re-solving $(A_0,E_4)$ on the fly for each first block, which makes the cost of the first block negligible; and
(iii)~a meet-in-the-middle lookup on 19~bits of $E_{16}$ whose table fits in the GPU's L2 cache.
The collision was found after about 3~hours of wall-clock time on three NVIDIA Blackwell GPUs (one RTX~5090 and two RTX~PRO~6000), i.e.\ about $8.9$ GPU-hours. We estimate this at roughly $2^{49}$ 36-step compression-function evaluations. The measured stage-by-stage pass rates agree with the cost model; for example, the step-22 survivors that also satisfy $\delta W_{23}$ number 9, against 9.3 expected. We verified the pair with an independent implementation, and independently re-measured the tail probabilities on $2\cdot 10^{10}$ generated message pairs.
\end{abstract}

\paragraph{AI-use designation.} Predominantly AI-generated work. This work was developed with AI agents under the author's direction. Google Gemini~4 Neon, running in the Antigravity IDE, wrote the code and ran the experiments. Anthropic Claude Opus~5.5 contributed the initial exploration, the independent verification in Sect.~\ref{sec:verify}, and the writing. The author takes responsibility for the results and presentation.

\section{Introduction}

SHA-512~\cite{FIPS180} has 80 steps. Collision attacks on step-reduced SHA-2 measure its security margin. Since 2024, automated searches for signed differential characteristics have combined with memory-efficient two-block frameworks, and the reach of these attacks has grown quickly~\cite{LLW24a,LLWDS24,LLWQJK25,LLWS26,LZLLQZ26}. Table~\ref{tab:sota} summarizes the collision attacks on SHA-512 with the standard IV, using the list in~\cite[Table~1]{LZLLQZ26}.

\begin{table}[h]
\centering
\caption{Collision attacks on step-reduced SHA-512 with the standard IV (after~\cite[Table~1]{LZLLQZ26}). Time is in calls to the compression function.}
\label{tab:sota}
\begin{tabular}{rlll}
\toprule
Steps & Time & Status & Reference \\
\midrule
27 & --- & practical & \cite{DEM15} \\
28 & --- & practical & \cite{LLW24a} \\
29 & --- & practical & \cite{LLWQJK25} \\
31 & $2^{85.5}$ & theoretical & \cite{LLWQJK25} \\
35 & $\approx 2^{54}$ & practical & \cite{LLWS26} \\
36 & $2^{136}$ & theoretical & \cite{LLWS26} \\
36 & $2^{66.2}$ & theoretical & \cite{LZLLQZ26} \\
\textbf{36} & $\approx 2^{49}$ (measured; see Sect.~\ref{sec:compare}) & \textbf{practical} & \textbf{this work} \\
37 & $2^{87}$ & theoretical & \cite{LZLLQZ26} \\
38 & $2^{125.4}$ & theoretical & \cite{LZLLQZ26} \\
39 & $2^{178}$ & theoretical & \cite{LZLLQZ26} \\
\bottomrule
\end{tabular}
\end{table}

For comparison, the 35-step practical collision for SHA-512 took about $2^{54}$ work and 14~days on one RTX~5060~Ti~\cite{LLWS26}. For SHA-256, Goldzweig recently gave a 38-step collision from the standard IV, found in 53~GPU-hours~\cite{Gol26}.

\paragraph{Contributions.}
\begin{itemize}
\item A colliding message pair for 36-step SHA-512 with the standard IV (Table~\ref{tab:pair}). It satisfies the standard padding as well: two 256-byte messages with equal 36-step SHA-512 digests.
\item A practical instantiation of the 36-step attack of~\cite{LZLLQZ26}. Three changes lower the expected cost from $2^{66.2}$ to about $2^{50}$ compression-function equivalents: a refined count of the degrees of freedom and conditions, on-the-fly re-solving of $(A_0,E_4)$ for each first block, and a meet-in-the-middle lookup resident in L2 (Sect.~\ref{sec:impl}).
\item Measured stage-wise statistics from $2.04\cdot 10^{9}$ step-18 survivors, and an independent audit of the probabilistic tail (Sect.~\ref{sec:verify}).
\end{itemize}
We stress that the differential characteristic, the message-difference pattern and the attack framework are those of Li~et~al.~\cite{LZLLQZ26,LLWS26}. We also reuse the internal-state solution of their published 36-step semi-free-start pair~\cite[Table~14]{LZLLQZ26} instead of solving that part again. None of the techniques in Sect.~\ref{sec:impl} is new in itself.

\section{Preliminaries}

SHA-512 processes 1024-bit blocks with 64-bit words. Let $\boxplus,\boxminus$ denote addition and subtraction modulo $2^{64}$, and let $\delta x = x'\boxminus x$ be the modular difference. For $16\le i<r$ the message expansion is
\[
W_i = \sigma_1(W_{i-2}) \boxplus W_{i-7} \boxplus \sigma_0(W_{i-15}) \boxplus W_{i-16},
\]
with $\sigma_0(x)=(x\ggg 1)\oplus(x\ggg 8)\oplus(x\gg 7)$ and $\sigma_1(x)=(x\ggg 19)\oplus(x\ggg 61)\oplus(x\gg 6)$. We use the usual two-register description, as in~\cite{LZLLQZ26}. Step $i$ computes
\begin{align*}
E_i &= A_{i-4} \boxplus E_{i-4} \boxplus \Sigma_1(E_{i-1}) \boxplus \mathrm{IF}(E_{i-1},E_{i-2},E_{i-3}) \boxplus K_i \boxplus W_i,\\
A_i &= E_i \boxminus A_{i-4} \boxplus \Sigma_0(A_{i-1}) \boxplus \mathrm{MAJ}(A_{i-1},A_{i-2},A_{i-3}),
\end{align*}
where $(A_{-4},\dots,A_{-1},E_{-4},\dots,E_{-1})$ is the chaining input. The $r$-step compression function adds the chaining input to the final state (feed-forward). A collision with the standard IV is a pair of messages $M_0\|M_1$ and $M_0\|M_1'$ with $M_1\neq M_1'$ that share the first block $M_0$ and give the same chaining value after the second block.

\section{The 36-step attack of Li et al.}\label{sec:prior}

The characteristic of~\cite[Tables~12--13]{LZLLQZ26} uses the message-difference family $i=1$ of~\cite{LLWS26}. It puts differences in $(W_5,\dots,W_9,W_{13},W_{14})$. After expansion, the only further non-zero differences among $W_{16},\dots,W_{35}$ are $\delta W_{21}=-(2^{57}+2^{51}+2^{43})$ and $\delta W_{23}=2^{59}$. The message differences realized by our pair are listed in Table~\ref{tab:pair}. The attack follows the improved memory-efficient framework of~\cite{LLWS26}:
\begin{enumerate}
\item Find a solution for $(W_i)_{8\le i\le 13}$, $(A_i)_{0\le i\le 13}$ and $(E_i)_{4\le i\le 13}$ that conforms to the middle part of the characteristic. The cost is about $2^{39.4}$ with a SAT solver.
\item For random first blocks $M_0$, compute the chaining value $\mathit{CV}_1$. Check the conditions on $(E_3,W_5,W_6,W_7)$, which match $(A_{-3},A_{-2},A_{-1})$; there are $4+4+6+4$ of them, and empirically a valid $M_0$ occurs with probability $2^{-18}$. Then derive $(W_0,\dots,W_4)$.
\item Use the freedom in $(W_{14},W_{15})$, estimated at $2^{21}$ values, to satisfy the 66 uncontrolled conditions on $(A_i,E_i,W_i)_{i\ge 16}$.
\end{enumerate}
This needs $2^{66-21}=2^{45}$ valid $M_0$. The total cost is estimated at $2^{66}+2^{45}\cdot 2^{18}\approx 2^{66.2}$~\cite{LZLLQZ26}.

\section{Making the attack practical}\label{sec:impl}

\subsection{Degrees of freedom in the second block}

We reuse the Step-1 solution embedded in the published semi-free-start pair~\cite[Table~14]{LZLLQZ26}. We fix $(A_i)_{1\le i\le 13}$, $(E_i)_{5\le i\le 13}$ and $(W_i)_{9\le i\le 13}$ to its values.

\emph{Upper part.} Four bits of $A_{14}$, at positions $\{28,40,62,63\}$, carry no condition. Changing $A_{14}$ by $\Delta$ changes $E_{14}$ and $W_{14}$ by the same $\Delta$, which gives 16 choices of $(A_{14},E_{14},W_{14})$. For each choice, a bit-level dynamic program over the carries of $E_{15}=A_{15}\boxplus(\cdot)$ enumerates the values of $A_{15}$ that meet all bit conditions on $A_{15}$ and $E_{15}$, plus three parity conditions. This yields exactly $2^{21}$ values of $A_{15}$, and hence of $W_{15}$, per choice: $2^{25}$ \emph{upper} candidates in total, against $2^{21}$ in~\cite{LZLLQZ26}. All 16 values of $W_{14}$ satisfy both $\sigma$-difference conditions on $W_{14}$ that are needed for $\delta W_{16}=0$ and $\delta W_{29}=0$; we checked this exhaustively.

\emph{Lower part: re-solving $(A_0,E_4)$ on the fly.} For a fixed $\mathit{CV}_1$ we keep $T_4:=E_4\boxminus A_0$ fixed and let $E_4$ range over its 29 unconditioned bits, at positions $\{0..3,5..17,20..24,27..29,34..36,63\}$. Setting $A_0=E_4\boxminus T_4$ leaves $A_4$ and every state word from step~5 on unchanged. The words $(E_0,\dots,E_3)$ and $(W_0,\dots,W_8)$ are recomputed by the formulas of~\cite[Step~2]{LZLLQZ26}. $W_8$ now depends on $E_4$, so its $\sigma_0$-difference condition must be checked again. This idea is close in spirit to the reuse of state freedom in~\cite[Sect.~4]{LLWS26}. Before the enumeration, a first block is accepted (``golden'') if the 14 conditions on $E_3$, $W_6$ and $W_7$ hold at the reference value of $E_4$. Over the $2^{29}$ choices of $E_4$, the conditions on $(E_3,W_5,W_6,W_7,W_8)$ then hold with measured probability $3.85\%$ ($2^{-4.70}$). That gives about $2^{24.3}$ \emph{lower} candidates per golden $M_0$, and the first block accounts for only 0.02\% of the running time. The four conditions on $W_5$ reduce to the bitwise test $W_5[59]=W_5[2]=0$, $W_5[52]=W_5[58]$, $W_5[53]\neq W_5[60]$. We checked on $1.3\cdot 10^8$ samples that this test agrees exactly with $\delta\sigma_0(W_5)=-\delta W_{13}$.

\subsection{Meet-in-the-middle on $E_{16}$}

Write $W_{16}=\sigma_1(W_{14})\boxplus W_9\boxplus\sigma_0(W_1)\boxplus W_0$. Then
\[
E_{16} = \ell \boxplus h_E, \qquad A_{16} = \ell \boxplus h_A, \qquad \ell := \sigma_0(W_1)\boxplus W_0,
\]
where $\ell$ depends only on the lower part, and $h_E,h_A$ depend only on the upper part. The bit conditions on $E_{16}$ and $A_{16}$ therefore turn into a matching problem on $\ell\boxplus h$. We sort the $2^{25}$ upper candidates by a 19-bit key on $h$, and file each lower candidate under the corresponding key on $\ell$. The key consists of 1~exact bit ($E_{16}[0]$); 7~bits from three carry-guarded slices ($E_{16}[15..16]$, $[17..18]$, $[36..38]$); 10~bits of $E_{16}$ predicted modulo carries; and $A_{16}[59]$. Of the $2^{19}$ buckets, 61{,}365 are non-empty. With at most 64 lower candidates per bucket, the hot table is about 30\,MB and stays in the L2 cache of the GPUs used (96\,MB on the RTX~5090). It captures 99.0\% of the lower candidates that hit a non-empty bucket.

\subsection{Filtering pipeline and cost model}

Every pair of a lower and an upper candidate in the same bucket goes through these filters in order:
\begin{enumerate}
\item 15 direct bits of $E_{16}$: 12 for step~17, plus 3 that step~18 forces ($E_{16}[43]=1$, $E_{16}[51]=E_{16}[57]=0$).
\item Two XOR conditions on $E_{16}$, and $A_{16}[59]=0$.
\item Compute $E_{17}$; require that the step-18 difference vanishes and that $A_{17}[59]=0$. Pairs that pass are called \emph{step-18 survivors}.
\item Compute both messages fully and check steps 18--35. The checks are $\delta A=\delta E=0$ where the characteristic requires it, $\delta E_{19}=-2^{59}$, then $\delta W_{23}=2^{59}$, then $\delta W_{25}=0$, and finally equality of the outputs.
\end{enumerate}
Once $\delta W_{23}$ and $\delta W_{25}$ are correct, no difference remains in the state or in $W_{24},\dots,W_{35}$, so equal outputs follow. We confirm this empirically in Sect.~\ref{sec:verify}. The expected number of step-18 survivors per collision is
\[
2^{16.19}\ (\text{steps }19\text{--}22)\;\cdot\;2^{11.51}\ (\delta W_{23})\;\cdot\;2^{3.83}\ (\delta W_{25},\ \text{probability }9/128) \;=\; 2^{31.53}.
\]
The two message-only factors were calibrated by Monte Carlo with $4\cdot 10^9$ samples each. Steps 19--22 were calibrated on sampled step-18 states. On the hardware below, one step-18 survivor costs $2^{21.98}$ bucket-matched pairs, so a collision is expected after $2^{53.5}$ pairs.

\section{Results}\label{sec:results}

\paragraph{Setup.} Three GPUs ran disjoint ranges of first-block seeds: one GeForce RTX~5090 (32\,GB) and two RTX~PRO~6000 Blackwell Server Edition. The code is a single CUDA program compiled for \texttt{sm\_120}. Measured throughput was 272.6\,G pairs/s and 65.8\,k step-18 survivors/s on the RTX~5090, and 253.5\,G pairs/s and 61.2\,k survivors/s on each RTX~PRO~6000. At this rate the model predicts 13.1~h on the RTX~5090 alone and 4.6~h on all three GPUs together.

\paragraph{Outcome.} The RTX~5090 found the collision after 10{,}113\,s (2.81\,h). By then the two other GPUs had run for about 10{,}950\,s each, so the total effort was about 8.9 GPU-hours. Table~\ref{tab:stats} lists the counters. Up to the collision, about $2.0\cdot 10^9$ step-18 survivors had been processed. That is 65\% of the expected $2^{31.53}$, close to the median of $\ln 2\approx 69\%$. Table~\ref{tab:model} compares the model with the observations.

\begin{table}[h]
\centering
\caption{Search counters. The ``$\delta W_{23}$'' and ``$\delta W_{25}$'' columns count step-22 survivors that also satisfy each message condition. The RTX~PRO~6000 runs continued for about 285\,s after the collision, until a stop file arrived; their counters include that time.}
\label{tab:stats}
\begin{tabular}{lrrrrrrr}
\toprule
GPU & time (s) & golden $M_0$ & step-18 & step-22 & $\delta W_{23}$ & $\delta W_{25}$ & collision \\
\midrule
RTX 5090 & 10{,}113 & 2{,}102{,}288 & 665{,}107{,}101 & 8{,}845 & 5 & 1 & 1 \\
RTX PRO 6000 \#1 & 11{,}240 & 2{,}173{,}696 & 687{,}390{,}638 & 9{,}131 & 1 & 0 & 0 \\
RTX PRO 6000 \#2 & 11{,}240 & 2{,}174{,}720 & 687{,}630{,}738 & 9{,}147 & 3 & 0 & 0 \\
\midrule
Total & 32{,}592 & 6{,}450{,}704 & 2{,}040{,}128{,}477 & 27{,}123 & 9 & 1 & 1 \\
\bottomrule
\end{tabular}
\end{table}

\begin{table}[h]
\centering
\caption{Predicted against observed stage probabilities (whole run).}
\label{tab:model}
\begin{tabular}{lll}
\toprule
Stage & Model & Observed \\
\midrule
step-18 $\to$ step-22 & $2^{-16.19}$ & $27{,}123/2.040\cdot10^9 = 2^{-16.20}$ \\
step-22 $\to$ $\delta W_{23}$ & $2^{-11.51}$ (9.3 expected) & $9/27{,}123 = 2^{-11.56}$ \\
$\delta W_{23}\to\delta W_{25}$ & $9/128=0.070$ (0.63 expected) & $1/9$ \\
$\delta W_{25}\to$ collision & $1$ & $1/1$ \\
\bottomrule
\end{tabular}
\end{table}

\begin{table}[t]
\centering
\caption{The colliding message pair for 36-step SHA-512 with the standard IV. Words are 64-bit, big-endian, listed in order. $M_0$ encodes the seed \hx{0x1e1b34abbd2} followed by the ASCII strings ``\hx{SHA51236}'' and ``\hx{COLLIDE!}''.}
\label{tab:pair}
\small
\setlength{\tabcolsep}{4pt}
\begin{tabular}{l@{\quad}llll}
\toprule
$M_0$ & \hx{000001e1b34abbd2} & \hx{5348413531323336} & \hx{434f4c4c49444521} & \hx{0000000000000000} \\
      & \multicolumn{4}{l}{(words 4--15 are \hx{0000000000000000})} \\
\midrule
$\mathit{CV}_1$ & \hx{f1e42e9f77e4cca4} & \hx{a785f886ea02658d} & \hx{cf52ba6f18e44e11} & \hx{644ba4c897d35d12} \\
                & \hx{c4cd7917344f6760} & \hx{3f554360df46b76a} & \hx{1a8281d8e746e746} & \hx{0ef673a2b4c69e41} \\
\midrule
$M_1$ & \hx{4f0a777b3e20699b} & \hx{9084d21a7d15ffdb} & \hx{f3eb37ed7c49a6a9} & \hx{811b8eea0f813426} \\
      & \hx{9901a1886b80a163} & \hx{36982650a4377b49} & \hx{e2734f1dfc05e77e} & \hx{65f45efd8b33bb64} \\
      & \hx{d395fbe4be6d579f} & \hx{b0f8774e493b43bc} & \hx{13287fa0a341d877} & \hx{f0d19e5ae32f6abf} \\
      & \hx{f26d943cf8041873} & \hx{06f2b17bb37eb2b7} & \hx{2cccf09ec39c66dd} & \hx{3284b0682fc4054b} \\
\midrule
$M_1'$ & \hx{4f0a777b3e20699b} & \hx{9084d21a7d15ffdb} & \hx{f3eb37ed7c49a6a9} & \hx{811b8eea0f813426} \\
       & \hx{9901a1886b80a163} & \hx{3e982650a4377b49} & \hx{e67b4f1dfc05e77e} & \hx{6df45efd8b33bb64} \\
       & \hx{f3c5db65be6d579f} & \hx{f0d8764e493b43bc} & \hx{13287fa0a341d877} & \hx{f0d19e5ae32f6abf} \\
       & \hx{f26d943cf8041873} & \hx{02fab17bb37eb2b7} & \hx{24ccf09ec39c66dd} & \hx{3284b0682fc4054b} \\
\midrule
$H$ & \hx{7f987d4f13fdfebe} & \hx{99203ca3a1bf5993} & \hx{1e4de96de6f5fbe5} & \hx{ba20aadb881fc9cf} \\
    & \hx{33724a7027ab508b} & \hx{3e25595a172fb3c8} & \hx{6b67ab3917239813} & \hx{209bb9516ff380cf} \\
\bottomrule
\end{tabular}

\medskip
\footnotesize
$\delta W_5=\hx{0800000000000000}$, $\delta W_6=\hx{0408000000000000}$, $\delta W_7=\hx{0800000000000000}$, $\delta W_8=\hx{202fdf8100000000}$,\\
$\delta W_9=\hx{3fdfff0000000000}$, $\delta W_{13}=\hx{fc08000000000000}$, $\delta W_{14}=\hx{f800000000000000}$; all other $\delta W_i$ ($i<16$) are zero.\\
With standard padding, the 256-byte messages $M_0\|M_1$ and $M_0\|M_1'$ have the same 36-step SHA-512 digest:\\
\hx{9eabd799baf71c41b426c3833361abcd9b9405a24d6ff6fa41f43e6071e437f3}\\
\hx{79413d838b9f4634be19868ac36dca673d3e031474cf1a596f05c4e4546e9e84}.
\end{table}

\section{Independent verification}\label{sec:verify}

\paragraph{The pair.} We wrote a second implementation of SHA-512 from scratch and derived its round constants and IV from the cube and square roots of the first primes. For the full 80 steps it agrees with Python's \texttt{hashlib} on test vectors. Truncated to 36 steps, it confirms the following:
\begin{itemize}
\item $\mathit{CV}_1$ is the compression of $M_0$ from the standard IV;
\item $M_1$ and $M_1'$ differ exactly in words $\{5,6,7,8,9,13,14\}$;
\item both messages compress to the same $H$;
\item the padded 256-byte messages collide as well, while their full 80-step SHA-512 digests differ.
\end{itemize}
The script is given in Appendix~\ref{app:verify}.

\paragraph{The tail model.} The $\delta W_{23}$ and $\delta W_{25}$ conditions depend only on the message words. Their probabilities could therefore be biased by the way the search constructs its candidates. They could also be correlated with the state conditions of steps 17--22. To test both, a separate CPU program generates candidate pairs exactly as the search does: golden first blocks, then random $E_4$ variants, then random upper candidates. Given message~1's real state and the characteristic's modular state difference before each step $t$, it checks whether step $t$ reproduces the characteristic's difference. Table~\ref{tab:audit} summarizes $2\cdot 10^{10}$ such pairs. The conditional probabilities of $\delta W_{23}$ given steps 17, 19, 20, 21 and 22 all lie within $2^{-11.50}$ to $2^{-11.56}$ (standard errors 0.001--0.10 bits). Very few random pairs satisfy the step-17/18 conditions, so the conditioning on those steps is only checked to about $\pm 0.1$ bit.

\begin{table}[h]
\centering
\caption{CPU audit on $2\cdot10^{10}$ generated candidate pairs.}
\label{tab:audit}
\begin{tabular}{lll}
\toprule
Quantity & Model & Measured \\
\midrule
$\Pr[\delta W_{23}=2^{59}]$ & $2^{-11.5114}$ & $2^{-11.5104\pm0.0006}$ \\
$\Pr[\delta W_{25}=0\mid \delta W_{23}=2^{59}]$ & $9/128=0.07031$ & $0.07041$ \\
steps 24--35 hold $\mid$ both & $1$ & $482{,}706/482{,}706$ \\
steps 0--16 hold for generated candidates & $1$ & always \\
$W_5$ bitwise test $\equiv$ $\sigma_0$ condition & exact & 0 mismatches in $1.28\cdot10^8$ \\
\bottomrule
\end{tabular}
\end{table}

\section{Comparison with the $2^{66.2}$ estimate}\label{sec:compare}

The estimate in~\cite{LZLLQZ26} counts calls to the compression function. The search up to the collision enumerated about $8.3\cdot 10^{15}\approx 2^{52.9}$ bucket-matched pairs. A pair check costs a handful of 64-bit integer operations, not a compression call. To convert, we assume a throughput of about $2.2\cdot 10^{10}$ 36-step SHA-512 compressions per second per GPU. We extrapolated this from public hashcat benchmarks of full SHA-512 on the RTX~5090 \todo{measure 36-step compression throughput directly on the RTX~5090 and replace this estimate}. On that basis, the 8.9 GPU-hours amount to roughly $2^{49}$ compression calls; the expected cost is about $2^{50}$. At the same throughput, $2^{66.2}$ compression calls would take about 39~years on the three GPUs. Table~\ref{tab:gap} gives an approximate attribution of the gap of about $2^{16}$--$2^{17}$.

\begin{table}[h]
\centering
\caption{Approximate attribution of the gap between $2^{66.2}$ and the expected cost of about $2^{50}$ (in bits).}
\label{tab:gap}
\begin{tabular}{lr}
\toprule
Source & bits \\
\midrule
Uncontrolled conditions recounted on real transitions ($66\to 62.2$) & $\approx 3.8$ \\
Bucket-matched enumeration: $2^{62.2}$ combinations covered by $2^{53.5}$ enumerated pairs & $\approx 8.7$ \\
Cost of one pair check, relative to one compression call (estimated) & $\approx 3.6$ \\
First-block term $2^{63}$ removed by on-the-fly re-solving of $(A_0,E_4)$ & $\approx 0.2$ \\
\midrule
Total & $\approx 16.3$ \\
\bottomrule
\end{tabular}
\end{table}

The largest share comes from the standard meet-in-the-middle matching. The original authors could have applied it themselves. The $2^{66.2}$ figure is best read as a conservative upper estimate, not as an error.

\section{Conclusion}

36-step SHA-512 with the standard IV is now practically broken: the collision took about 3~hours on three commodity Blackwell GPUs. The measured pass rates agree with the cost model to within statistical error. Full SHA-512 has 80 steps, and this result has no implication for its security. The next steps of the same family, 37~steps at $2^{87}$ and 38~steps at $2^{125.4}$~\cite{LZLLQZ26}, are outside the reach of the techniques used here; we have not tested whether similar refinements reduce those estimates.

\paragraph{Data availability.} \todo{Public repository URL with the CUDA search code, the CPU audit, and the verification script.}

\begin{thebibliography}{99}
\bibitem{DEM15} C.~Dobraunig, M.~Eichlseder, F.~Mendel. Analysis of SHA-512/224 and SHA-512/256. ASIACRYPT 2015, LNCS, pp.~612--630. Springer, 2015.
\bibitem{FIPS180} National Institute of Standards and Technology. FIPS 180-4: Secure Hash Standard (SHS), 2012.
\bibitem{Gol26} O.~Goldzweig. A collision for 38-step SHA-256. Cryptology ePrint Archive, Paper 2026/2913, 2026.
\bibitem{LLW24a} Y.~Li, F.~Liu, G.~Wang. New records in collision attacks on SHA-2. EUROCRYPT 2024, LNCS, pp.~158--186. Springer, 2024.
\bibitem{LLWDS24} Y.~Li, F.~Liu, G.~Wang, X.~Dong, S.~Sun. The first practical collision for 31-step SHA-256. ASIACRYPT 2024, LNCS, pp.~237--266. Springer, 2024.
\bibitem{LLWQJK25} Y.~Li, F.~Liu, G.~Wang, H.~Qian, K.~Jia, X.~Kong. New collision attacks on round-reduced SHA-512. CRYPTO 2025, LNCS, pp.~200--229. Springer, 2025.
\bibitem{LLWS26} Y.~Li, F.~Liu, G.~Wang, J.~Shi. Pushing the limit of memory-efficient collision attack framework for SHA-2. CRYPTO 2026. Cryptology ePrint Archive, Paper 2026/1080, 2026.
\bibitem{LZLLQZ26} Y.~Li, Z.~Zhang, M.~Li, F.~Liu, H.~Qian, J.~Zhu. Pushing collision attacks on SHA-2 to 39 steps. Cryptology ePrint Archive, Paper 2026/1120, 2026.
\end{thebibliography}

\appendix
\section{Verification script}\label{app:verify}
The script below needs only the Python standard library. It reads the record of the pair in JSON form.
\begin{lstlisting}[language=Python]
import json, sys, hashlib, struct
M64 = (1 << 64) - 1
def primes(n):
    ps, x = [], 2
    while len(ps) < n:
        if all(x % p for p in ps if p * p <= x): ps.append(x)
        x += 1
    return ps
def iroot(n, k):
    lo, hi = 0, 1 << (n.bit_length() // k + 1)
    while lo < hi:
        mid = (lo + hi + 1) // 2
        lo, hi = (mid, hi) if mid ** k <= n else (lo, mid - 1)
    return lo
P = primes(80)
K  = [iroot(p << 192, 3) & M64 for p in P]      # frac(cbrt p) * 2^64
IV = [iroot(p << 128, 2) & M64 for p in P[:8]]  # frac(sqrt p) * 2^64
rotr = lambda x, n: ((x >> n) | (x << (64 - n))) & M64
S0 = lambda x: rotr(x, 28) ^ rotr(x, 34) ^ rotr(x, 39)
S1 = lambda x: rotr(x, 14) ^ rotr(x, 18) ^ rotr(x, 41)
s0 = lambda x: rotr(x, 1) ^ rotr(x, 8) ^ (x >> 7)
s1 = lambda x: rotr(x, 19) ^ rotr(x, 61) ^ (x >> 6)
def compress(cv, M, steps):
    W = list(M) + [0] * (steps - 16)
    for i in range(16, steps):
        W[i] = (s1(W[i-2]) + W[i-7] + s0(W[i-15]) + W[i-16]) & M64
    a, b, c, d, e, f, g, h = cv
    for t in range(steps):
        t1 = (h + S1(e) + ((e & f) ^ (~e & g)) + K[t] + W[t]) & M64
        t2 = (S0(a) + ((a & b) ^ (a & c) ^ (b & c))) & M64
        h, g, f, e, d, c, b, a = g, f, e, (d + t1) & M64, c, b, a, (t1 + t2) & M64
    return [(x + y) & M64 for x, y in zip(cv, [a, b, c, d, e, f, g, h])]
def blocks(m):
    L = len(m) * 8
    m += b"\x80" + b"\x00" * ((112 - (len(m) + 1) % 128) % 128) + L.to_bytes(16, "big")
    return [list(struct.unpack(">16Q", m[i:i+128])) for i in range(0, len(m), 128)]
def sha512r(m, steps):
    cv = IV
    for B in blocks(m): cv = compress(cv, B, steps)
    return b"".join(x.to_bytes(8, "big") for x in cv)
for m in [b"", b"abc", bytes(range(256)) * 3]:
    assert sha512r(m, 80) == hashlib.sha512(m).digest()
J = json.load(open(sys.argv[1])); h = lambda L: [int(x, 16) for x in L]
M0, M1, M2 = h(J["M0"]), h(J["M1"]), h(J["M2"])
cv = compress(IV, M0, 36)
assert M1 != M2 and compress(cv, M1, 36) == compress(cv, M2, 36)
a = b"".join(x.to_bytes(8, "big") for x in M0 + M1)
b = b"".join(x.to_bytes(8, "big") for x in M0 + M2)
assert sha512r(a, 36) == sha512r(b, 36) and hashlib.sha512(a).digest() != hashlib.sha512(b).digest()
print("36-step SHA-512 collision verified:", sha512r(a, 36).hex())
\end{lstlisting}

\end{document}
